\chapter{Evaluation}
\label{ch:evaluation}

Die Evaluation soll zeigen, ob das Gesamtsystem die Zielsetzung aus \cref{ch:einleitung}
erfüllt: Ein Bauteil wird ohne Eingriff eines Bedieners erkannt, gegriffen, geprüft und
richtig sortiert. Sie ist als Funktionsnachweis angelegt und nicht als Messkampagne zur
Genauigkeit einzelner Verfahren. Grundlage sind die Zyklusprotokolle aus
\cref{sec:int-protokoll}.

\section{Versuchsaufbau und Vorgehen}
\label{sec:eval-vorgehen}

Alle Versuche finden an der realen Anlage statt, mit allen Modulen in der Betriebsstufe
\enquote{Real} (\cref{sec:teststrategie}). Es werden zwei Versuchsreihen durchgeführt.

\paragraph{Versuchsreihe 1: Wiederholbarkeit.}
Ein Gutteil wird wiederholt von Hand an zufälliger Position und mit zufälliger Orientierung in
den Arbeitsbereich gelegt, und jeweils ein vollständiger Zyklus wird gestartet. Die Reihe zeigt,
ob Erkennung, Greifen, Ablegen an der Prüfposition und Prüfung für unterschiedliche Lagen
zuverlässig zusammenspielen und das Teil jedes Mal als Gutteil sortiert wird.

\paragraph{Versuchsreihe 2: Aussortierung.}
Es werden gezielt fehlerhafte Bauteile gefertigt, bei denen jeweils eines der Prüfmerkmale
nicht erfüllt ist (\cref{tab:eval-fehlerteile}). Jedes dieser Teile durchläuft den Zyklus wie
in Versuchsreihe~1. Die Reihe zeigt, ob das System die Fehler erkennt und die Teile als
Schlechtteil ablegt.

\begin{table}[htbp]
  \centering
  \caption{Fehlerhafte Bauteile für Versuchsreihe 2}
  \label{tab:eval-fehlerteile}
  \begin{tabularx}{\textwidth}{@{}l X X@{}}
    \toprule
    Teil & Fehler & Erwartetes Ergebnis \\
    \midrule
    A & Kerbe fehlt & Schlechtteil, Grund \enquote{Kerbe fehlt} \\
    B & Lochdurchmesser außerhalb der Toleranz & Schlechtteil, Grund \enquote{Durchmesser außerhalb der Toleranz} \\
    C & Etikett mit Seriennummer entfernt & Schlechtteil, Grund \enquote{Seriennummer ungültig} \\
    \bottomrule
  \end{tabularx}
\end{table}

\todo[inline]{Anzahl der Durchläufe je Versuchsreihe festlegen. Festlegen, ob und welche
Kennzahlen ausgewertet werden (z.\,B.\ Anteil erfolgreicher Zyklen, Fehler je Zustand,
Zyklusdauer).}

\section{Ergebnisse Versuchsreihe 1: Wiederholbarkeit}
\label{sec:eval-vr1}
\todo[inline]{Ergebnisse nach den Versuchen ergänzen.}

\section{Ergebnisse Versuchsreihe 2: Aussortierung}
\label{sec:eval-vr2}
\todo[inline]{Ergebnisse nach den Versuchen ergänzen.}

\section{Vergleich klassischer und KI-basierter Ansatz}
\label{sec:eval-vergleich}
\todo[inline]{Person 4: Vergleich der Objekterkennung (\cref{ch:objekterkennung}) mit dem
lernbasierten Ansatz (\cref{ch:kigreifen}).}
